% =========================================================================
% PLATFORM TAXONOMY ROUTING (patho.dentnet.space)
% =========================================================================
% DOC_ID: odontogenic_cysts
% DOC_SLUG: odontogenic_cysts
% DOC_CATEGORY: cat_odontogenic_cysts
% DOC_TITLE: الگوریتم جامع: کیست‌های ادنتوژنیک فکین (تشخیص افتراقی)
% DOC_ALIASES: odontogenic cysts, کیست‌های ادنتوژنیک, کیست رادیکولار, دنتیژروس, کراتوسیست, OKC, OOC, کیست رویشی
% =========================================================================
% === CLINICAL EVIDENCE LEDGER & PMIDS ===
% 1. WHO Classification of Head and Neck Tumours (5th Ed, 2024) - Odontogenic Cysts.
% 2. Neville BW, Damm DD, Allen CM, Chi AC. Oral and Maxillofacial Pathology (5th Ed).
% 3. PMID: 41454329 - Mahmoud R et al. (BMC Oral Health 2025) 5-FU Meta-Analysis in OKC.
% 4. PMID: 40269316 - Aarthi S et al. (Evid Based Dent 2025) Alternatives to Carnoy's Solution.
% 5. PMID: 39038594 - Jacobs T et al. (J Oral Maxillofac Surg 2024) 5-FU vs MCS Recurrence & Nerve Safety.
% 6. PMID: 39406309 - Cobo-Vázquez C et al. (J Stomatol Oral Maxillofac Surg 2025) Pediatric Dentigerous Cyst.
% 7. PMID: 42255106 - Leena Sankari S et al. (Case Rep Dent 2026) Orthokeratinized Odontogenic Cyst Criteria.
% 8. PMID: 42054731 - Qu J et al. (Arch Oral Biol 2026) SPARC and MMP-9 in OKC vs OOC.
% =========================================================================

\documentclass[a4paper]{article}
\usepackage[margin=1.15cm,footskip=0pt,headsep=0pt]{geometry}
\usepackage{amsmath,amssymb}
\usepackage{xcolor}
\usepackage{tikz}
\usetikzlibrary{arrows.meta,positioning,shapes.geometric}
\usepackage{hyperref}
\usepackage{fontspec}

\makeatletter
\newcommand{\MathCollectFalse}{}
\newcommand{\MathCollectTrue}{}
\let\@kernel@tabular@init\relax
\makeatother

\usepackage{xepersian}

\settextfont[
  Path = Peyda/,
  BoldFont = Peyda-Bold.ttf,
  ItalicFont = Peyda-Medium.ttf,
  BoldItalicFont = Peyda-Bold.ttf
]{Peyda-Regular.ttf}

\setlatintextfont{Times New Roman}

\IfFileExists{Peyda/Symbola.ttf}{%
  \newfontfamily\symb[Path=Peyda/]{Symbola.ttf}%
}{%
  \IfFileExists{editor/Peyda/Symbola.ttf}{%
    \newfontfamily\symb[Path=editor/Peyda/]{Symbola.ttf}%
  }{%
    \newfontfamily\symb{Symbola}%
  }%
}
\newcommand{\emoji}[1]{{\symb #1}}

\linespread{1.10}
\pagestyle{empty}
\thispagestyle{empty}

\raggedbottom
\setlength{\parindent}{0pt}
\setlength{\parskip}{0pt}

\definecolor{hlgray}{HTML}{E2E8F0}
\definecolor{hlyellow}{HTML}{FEF08A}
\definecolor{hlred}{HTML}{FEE2E2}
\definecolor{darkred}{HTML}{991B1B}
\definecolor{navy}{HTML}{0A2540}
\definecolor{teal}{HTML}{0D9488}

\begin{document}

\begin{center}
{\Large\bfseries فلوچارت تصمیم‌گیری بالینی: کیست‌های ادنتوژنیک فکین}\\[2pt]
{\normalsize\bfseries\color{navy} \lr{Odontogenic Cysts of the Jaws (Differential Management)}}\\[3pt]
{\small\itshape تریاژ حیات پالپ، آسپیراسیون مایع کیستیک، پایش کراتینیزاسیون دیواره، پروتکل‌های کارنوی/۵-FU و حفظ دندان}
\end{center}
\vspace{-4pt}
\hrule height 0.8pt
\vspace{6pt}

\begin{center}
\begin{latin}
\linespread{0.95}\selectfont
\begin{tikzpicture}[
    every node/.style={font=\sffamily\footnotesize, align=center},
    stepbox/.style={rectangle, draw=black, fill=white, line width=0.7pt, inner sep=4.5pt},
    decision/.style={diamond, aspect=2.4, draw=black, fill=white, line width=0.7pt, inner sep=2pt, font=\sffamily\scriptsize\bfseries},
    arrow/.style={-{Stealth[scale=0.9]}, thick, black}
]

% 1. Initial Presentation
\node (n1) [stepbox, line width=1.0pt, text width=12.2cm] at (0, 0) {
    \rl{\textbf{تظاهر بالینی و رادیوگرافیک: ضایعه رادیولوسنت کورتیکه‌دار فک (\lr{Jaw Radiolucency})}}\\[2pt]
    \scriptsize \rl{لوسنسی تک‌حفره‌ای یا چندحفره‌ای با مارجین مشخص $\bullet$ همراه با دندان نرویتال، نهفته یا بدون دندان $\bullet$ رد تروما و فیبروز}
};

% 2. Pulp Vitality & Aspiration Triage
\node (n2) [stepbox, text width=14.2cm] at (0, -1.90) {
    \rl{\colorbox{hlyellow}{\strut\textbf{مرحله ۱ الزامی: تست حیات پالپ دندان‌های مجاور، سی‌بی‌سی‌تی و آسپیراسیون سوزنی (\lr{FNA})}}}\\[2pt]
    \scriptsize \rl{سنجش حیات الکتریکی/حرارتی پالپ $\bullet$ رنگ مایع: کاهی (سروز)، چرکی، شیری/کرمی غلیظ (کراتین) $\bullet$ بررسی گسترش قدامی-خلفی}
};

% 3. Pulp Vitality Decision
\node (d1) [decision] at (0, -4.00) {\rl{حیات پالپ؟}\\[-1pt]\scriptsize \lr{Pulp Vitality?}};
\node (r_benign) [stepbox, draw=teal, line width=0.9pt, text width=3.8cm] at (5.8, -4.00) {
    \rl{\textbf{ضایعه پری‌آپیکال دندانی \emoji{🔍}}}\\[2pt]
    \scriptsize \rl{دندان نرویتال با ضایعه کوچک؛}\\\rl{درمان ریشه (\lr{RCT}) خط اول؛}\\\rl{بهبودی در ۷۰٪ تا ۸۰٪ بدون جراحی}
};

% 4. Mandatory Incisional Biopsy & Keratin Analysis
\node (n3) [stepbox, line width=1.0pt, text width=15.8cm] at (0, -6.40) {
    \rl{\colorbox{hlyellow}{\strut\textbf{بیوپسی برشی از دیواره کیست (\lr{Incisional Biopsy}) و ارزیابی کراتینیزاسیون و غربالگری سندرمیک \emoji{🔬}}}}\\[2pt]
    \scriptsize \rl{تفکیک پاراکراتینیزه در برابر ارتوکراتینیزه و غیرکراتینیزه $\bullet$ سنجش ضخامت اپیتلیوم دیواره $\bullet$ بررسی علائم بالینی سندرم گورلین (\lr{PTCH1})}
};

% 5. Subtype Decision
\node (d2) [decision] at (0, -8.70) {\rl{پاتولوژی دیواره}\\[-1pt]\scriptsize \lr{Histopathologic Subtype?}};

% 6. 4 Main Columns
\node (h1) [stepbox, text width=3.6cm] at (-6.15, -11.00) {
    \rl{\textbf{کیست رادیکولار مقاوم}}\\[2pt]
    \scriptsize \lr{Radicular / Inflammatory}\\\rl{اپیتلیوم مطبق غیرکراتینیزه}\\\rl{اجسام راشتون و ارتشاح التهابی}\\\rl{عود با جراحی کامل: \colorbox{hlgray}{\strut\textbf{\rl{کمتر از ۲٪}}}}
};
\node (h2) [stepbox, text width=3.6cm] at (-2.05, -11.00) {
    \rl{\textbf{کیست دنتی‌ژروس تکاملی}}\\[2pt]
    \scriptsize \lr{Dentigerous Cyst}\\\rl{اتصال به طوق (\lr{CEJ}) دندان نهفته}\\\rl{اپیتلیوم نازک ۲ تا ۴ لایه‌ای}\\\rl{عود با انوکلئیشن: \colorbox{hlgray}{\strut\textbf{\rl{کمتر از ۲٪}}}}
};
\node (h3) [stepbox, draw=darkred, line width=1.2pt, text width=3.6cm] at (2.05, -11.00) {
    \rl{\textbf{کراتوسیست ادنتوژنیک (\lr{OKC}) \emoji{⚡}}}\\[2pt]
    \scriptsize \lr{Odontogenic Keratocyst}\\\rl{پاراکراتین چین‌دار؛ بازال پالیسادی}\\\rl{موتاسیون \lr{PTCH1} / سندرمیک}\\\rl{عود با کورتاژ ساده: \colorbox{hlred}{\strut\color{darkred}\textbf{۴۰٪ تا ۶۰٪}}}
};
\node (h4) [stepbox, line width=0.9pt, text width=3.6cm] at (6.15, -11.00) {
    \rl{\textbf{کیست ارتوکراتینیزه / رویشی}}\\[2pt]
    \scriptsize \lr{Orthokeratinized / Eruption}\\\rl{ارتوکراتین ضخیم + لایه گرانولار}\\\rl{فاقد تهاجم مدولار و جهش ژنتیکی}\\\rl{عود با انوکلئیشن: \colorbox{hlgray}{\strut\textbf{\rl{کمتر از ۲٫۲٪}}}}
};

% 7. Primary Interventions
\node (i1) [stepbox, text width=3.6cm] at (-6.15, -13.50) {
    \rl{\textbf{انوکلئیشن + آپیکواکتومی}}\\[2pt]
    \scriptsize \rl{خارج‌سازی کامل دیواره کیست}\\\rl{قطع ۳ میلی‌متر انتهای ریشه}\\\rl{سیل رتروگراد با بیوسرامیک (\lr{MTA})}
};
\node (i2) [stepbox, text width=3.6cm] at (-2.05, -13.50) {
    \rl{\textbf{مارسوپیالیزاسیون یا انوکلئیشن}}\\[2pt]
    \scriptsize \textbf{\rl{دکمپرسیون در اطفال برای حفظ دندان}}\\\rl{\lr{[\textbf{PMID: 39406309}]}}\\\rl{یا انوکلئیشن + خارج‌سازی دندان}
};
\node (i3) [stepbox, draw=darkred, line width=1.2pt, text width=3.6cm] at (2.05, -13.50) {
    \rl{\textbf{انوکلئیشن + ادجوانت شیمیایی}}\\[2pt]
    \scriptsize \textbf{\rl{استکتومی محیطی + ۵-فلورواوراسیل}}\\\rl{گاز آغشته به \lr{5\% 5-FU} به مدت ۲۴ ساعت}\\\rl{کاهش عود به کمتر از ۸٪}\\\rl{\lr{[\textbf{PMID: 41454329}]}}
};
\node (i4) [stepbox, line width=0.9pt, text width=3.6cm] at (6.15, -13.50) {
    \rl{\textbf{انوکلئیشن محافظه‌کارانه}}\\[2pt]
    \scriptsize \rl{جداسازی تمیز کپسول از استخوان}\\\rl{پرهیز از کاربرد ادجوانت شیمیایی}\\\rl{اکسیزیون سقف در کیست رویشی}
};

% 8. Recall Schedules
\node (s1) [stepbox, text width=3.6cm] at (-6.15, -16.00) {
    \rl{\textbf{زمان‌بندی پیگیری:}}\\[2pt]
    \colorbox{hlgray}{\strut\textbf{\rl{سالانه تا ۲ سال}}}\\[2pt]
    \scriptsize \rl{پایش رادیوگرافیک پری‌آپیکال؛}\\\rl{تایید بازسازی استخوان کورتیکال}
};
\node (s2) [stepbox, text width=3.6cm] at (-2.05, -16.00) {
    \rl{\textbf{زمان‌بندی پیگیری:}}\\[2pt]
    \colorbox{hlgray}{\strut\textbf{\rl{هر ۶ ماه تا ۲ سال}}}\\[2pt]
    \scriptsize \rl{پایش مسیر رویش دندان دائم؛}\\\rl{انوکلئیشن تکمیلی در صورت نیاز}
};
\node (s3) [stepbox, draw=darkred, line width=1.2pt, text width=3.6cm] at (2.05, -16.00) {
    \rl{\textbf{زمان‌بندی پیگیری:}}\\[2pt]
    \colorbox{hlred}{\strut\color{darkred}\textbf{\rl{هر ۶ ماه (۵ سال)؛ سپس سالانه}}}\\[2pt]
    \scriptsize \rl{تصویربرداری \lr{CBCT} دوره‌ای؛}\\\rl{پایش کیست‌های اقماری ثانویه}
};
\node (s4) [stepbox, line width=0.9pt, text width=3.6cm] at (6.15, -16.00) {
    \rl{\textbf{زمان‌بندی پیگیری:}}\\[2pt]
    \colorbox{hlgray}{\strut\textbf{\rl{سالانه تا ۲ سال}}}\\[2pt]
    \scriptsize \rl{معاینه بالینی و پانورامیک ساده؛}\\\rl{پیش‌آگهی عالی و بدون عود}
};

% 9. Trigger Box
\node (trigger) [stepbox, line width=0.9pt, text width=16.4cm, align=right] at (0, -18.70) {
    \centerline{\rl{\textbf{هشدارهای بالینی و اندیکاسیون‌های تغییر پروتکل جراحی یا بیوپسی مجدد (\lr{Clinical Red Flags}):}}}\par\vspace{3pt}
    \scriptsize \rl{\colorbox{hlyellow}{\strut\textbf{مداخله فوری و ارتقای جراحی در:}} (۱) پرفوراسیون کورتکس باکال/لینگوال یا اتصال به بافت نرم پیرامونی $\bullet$ (۲) بی‌حسی لب یا چانه (\lr{Parethesia}) ناشی از فشار بر کانال عصب دندانی تحتانی $\bullet$ (۳) عود سریع رادیولوسنسی در کمتر از ۱۲ ماه پس از جراحی $\bullet$ (۴) وجود ندول سالید داخل لومن کیستیک (مشکوک به آملوبلاستوما یا دگرگونی بدخیم) $\bullet$ (۵) بروز کیست‌های چندگانه در فکین در سنین نوجوانی (معیار ماژور سندرم گورلین-گولتز).}
};

% 10. Refractory Box
\node (refract) [stepbox, text width=16.4cm, align=right] at (0, -21.20) {
    \centerline{\rl{\textbf{پروتکل موارد مقاوم، درمان‌های مولکولی و پایش کارسینوم اولیه داخل استخوان (\lr{PIOC}):}}}\par\vspace{3pt}
    \scriptsize \rl{\textbf{سطح ۱ (سالواژ در عودهای مکرر OKC):} رزکسیون مارجینال استخوان فک با حفظ لبه تحتانی مندیبل یا رزکسیون سگمنتال همراه با بازسازی فوری فلپ استخوانی در صورت نفوذ تهاجمی به بافت نرم خارج کورتکس.}\par\vspace{2.5pt}
    \scriptsize \rl{\textbf{سطح ۲ (مهار هدفمند مسیر مولکولی):} کاربرد تارگت‌تراپی مهارکننده مسیر \lr{Hedgehog} (\lr{Vismodegib}) در بیماران مبتلا به سندرم گورلین با کیست‌های راجعه و وسیع غیرقابل جراحی.}\par\vspace{2.5pt}
    \scriptsize \rl{\textbf{سطح ۳ (پایش بدخیمی):} بیوپسی کامل از تمام نواحی ضخیم‌شدگی دیواره کیست جهت رد کارسینوم سلول سنگفرشی اولیه با منشأ کیستیک (\lr{PIOSCC})؛ پرهیز اکید از باقی گذاشتن اپیتلیوم دیواره بدون هیستوپاتولوژی قطعی.}
};

% Connections
\draw [arrow] (n1) -- (n2);
\draw [arrow] (n2) -- (d1);
\draw [arrow] (d1.east) -- node[above, font=\scriptsize\bfseries] {\rl{نرویتال (پری‌آپیکال)}} (r_benign.west);
\draw [arrow] (d1.south) -- node[right=5pt, font=\scriptsize\bfseries] {\rl{ویتال / کیست وسیع}} (n3.north);
\draw [arrow] (n3.south) -- (d2.north);

% Bus-Bar from Decision 2 to 4 Subtype columns
\draw [thick, black] (d2.south) -- (0, -9.80);
\draw [thick, black] (-6.15, -9.80) -- (6.15, -9.80);
\draw [arrow] (-6.15, -9.80) -- (h1.north);
\draw [arrow] (-2.05, -9.80) -- (h2.north);
\draw [arrow] (2.05, -9.80) -- (h3.north);
\draw [arrow] (6.15, -9.80) -- (h4.north);

% Column Downward Arrows
\draw [arrow] (h1.south) -- (i1.north);
\draw [arrow] (h2.south) -- (i2.north);
\draw [arrow] (h3.south) -- (i3.north);
\draw [arrow] (h4.south) -- (i4.north);

\draw [arrow] (i1.south) -- (s1.north);
\draw [arrow] (i2.south) -- (s2.north);
\draw [arrow] (i3.south) -- (s3.north);
\draw [arrow] (i4.south) -- (s4.north);

% Collector Bus into Trigger Box
\draw [thick, black] (s1.south) -- (-6.15, -17.30);
\draw [thick, black] (s2.south) -- (-2.05, -17.30);
\draw [thick, black] (s3.south) -- (2.05, -17.30);
\draw [thick, black] (s4.south) -- (6.15, -17.30);
\draw [thick, black] (-6.15, -17.30) -- (6.15, -17.30);
\draw [arrow] (0, -17.30) -- (trigger.north);

% Final link to Refractory Box
\draw [arrow] (trigger.south) -- (refract.north);

\end{tikzpicture}
\end{latin}
\end{center}

\end{document}
